\documentclass{openreader}
\title{Ἰωάννου αʹ}
\date{}
\setmainlanguage{english}
\setmainfont{Charis SIL}
\setotherlanguage{greek}
\newfontfamily\greekfont[Script=Greek]{SBL BibLit}
\begin{document}
\ORBselectlanguage{greek}
\maketitle
\raggedbottom 
\fontsize{16pt}{24pt}\selectfont

\ChapterHeader{1}{Chapter 1}

\MainTextVerseMark{1}{1}Ὃ ἦν\FnParse{a}{imperfect active indicative 3rd singular} ἀπ’ ἀρχῆς, ὃ ἀκηκόαμεν,\FnParse{b}{perfect active indicative 1st plural} ὃ ἑωράκαμεν\FnParse{c}{perfect active indicative 1st plural} τοῖς ὀφθαλμοῖς ἡμῶν, ὃ ἐθεασάμεθα\FnParseFormGloss{d}{aorist middle indicative 1st plural}{θεάομαι}{to see} καὶ αἱ χεῖρες ἡμῶν ἐψηλάφησαν,\FnParseFormGloss{e}{aorist active indicative 3rd plural}{ψηλαφάω}{to touch} περὶ τοῦ λόγου τῆς ζωῆς— 
\MainTextVerseMark{1}{2}καὶ ἡ ζωὴ ἐφανερώθη,\FnParse{a}{aorist passive indicative 3rd singular} καὶ ἑωράκαμεν\FnParse{b}{perfect active indicative 1st plural} καὶ μαρτυροῦμεν\FnParse{c}{present active indicative 1st plural} καὶ ἀπαγγέλλομεν\FnParse{d}{present active indicative 1st plural} ὑμῖν τὴν ζωὴν τὴν αἰώνιον ἥτις ἦν\FnParse{e}{imperfect active indicative 3rd singular} πρὸς τὸν πατέρα καὶ ἐφανερώθη\FnParse{f}{aorist passive indicative 3rd singular} ἡμῖν— 
\MainTextVerseMark{1}{3}ὃ ἑωράκαμεν\FnParse{a}{perfect active indicative 1st plural} καὶ ἀκηκόαμεν\FnParse{b}{perfect active indicative 1st plural} ἀπαγγέλλομεν\FnParse{c}{present active indicative 1st plural} ⸀καὶ ὑμῖν, ἵνα καὶ ὑμεῖς κοινωνίαν\FnParseFormGloss{d}{accusative feminine singular}{κοινωνία}{fellowship} ἔχητε\FnParse{e}{present active subjunctive 2nd plural} μεθ’ ἡμῶν· καὶ ἡ κοινωνία\FnParseFormGloss{f}{nominative feminine singular}{κοινωνία}{fellowship} δὲ ἡ ἡμετέρα\FnParseFormGloss{g}{nominative feminine singular}{ἡμέτερος}{our} μετὰ τοῦ πατρὸς καὶ μετὰ τοῦ υἱοῦ αὐτοῦ Ἰησοῦ Χριστοῦ· 
\MainTextVerseMark{1}{4}καὶ ταῦτα γράφομεν\FnParse{a}{present active indicative 1st plural} ⸀ἡμεῖς ἵνα ἡ χαρὰ ἡμῶν ᾖ\FnParse{b}{present active subjunctive 3rd singular} πεπληρωμένη.\FnParse{c}{perfect passive participle nominative feminine singular} 
\MainTextVerseMark{1}{5}Καὶ ἔστιν\FnParse{a}{present active indicative 3rd singular} αὕτη ἡ ἀγγελία\FnParseFormGloss{b}{nominative feminine singular}{ἀγγελία}{message} ἣν ἀκηκόαμεν\FnParse{c}{perfect active indicative 1st plural} ἀπ’ αὐτοῦ καὶ ἀναγγέλλομεν\FnParseFormGloss{d}{present active indicative 1st plural}{ἀναγγέλλω}{to tell} ὑμῖν, ὅτι ὁ θεὸς φῶς ἐστιν\FnParse{e}{present active indicative 3rd singular} καὶ σκοτία\FnParseFormGloss{f}{nominative feminine singular}{σκοτία}{darkness} ⸂ἐν αὐτῷ οὐκ ἔστιν⸃\FnParse{g}{present active indicative 3rd singular} οὐδεμία. 
\MainTextVerseMark{1}{6}ἐὰν εἴπωμεν\FnParse{a}{aorist active subjunctive 1st plural} ὅτι κοινωνίαν\FnParseFormGloss{b}{accusative feminine singular}{κοινωνία}{fellowship} ἔχομεν\FnParse{c}{present active indicative 1st plural} μετ’ αὐτοῦ καὶ ἐν τῷ σκότει περιπατῶμεν,\FnParse{d}{present active subjunctive 1st plural} ψευδόμεθα\FnParseFormGloss{e}{present middle indicative 1st plural}{ψεύδομαι}{to lie} καὶ οὐ ποιοῦμεν\FnParse{f}{present active indicative 1st plural} τὴν ἀλήθειαν· 
\MainTextVerseMark{1}{7}ἐὰν δὲ ἐν τῷ φωτὶ περιπατῶμεν\FnParse{a}{present active subjunctive 1st plural} ὡς αὐτός ἐστιν\FnParse{b}{present active indicative 3rd singular} ἐν τῷ φωτί, κοινωνίαν\FnParseFormGloss{c}{accusative feminine singular}{κοινωνία}{fellowship} ἔχομεν\FnParse{d}{present active indicative 1st plural} μετ’ ἀλλήλων καὶ τὸ αἷμα ⸀Ἰησοῦ τοῦ υἱοῦ αὐτοῦ καθαρίζει\FnParse{e}{present active indicative 3rd singular} ἡμᾶς ἀπὸ πάσης ἁμαρτίας. 
\MainTextVerseMark{1}{8}ἐὰν εἴπωμεν\FnParse{a}{aorist active subjunctive 1st plural} ὅτι ἁμαρτίαν οὐκ ἔχομεν,\FnParse{b}{present active indicative 1st plural} ἑαυτοὺς πλανῶμεν\FnParse{c}{present active indicative 1st plural} καὶ ἡ ἀλήθεια ⸂οὐκ ἔστιν\FnParse{d}{present active indicative 3rd singular} ἐν ἡμῖν⸃. 
\MainTextVerseMark{1}{9}ἐὰν ὁμολογῶμεν\FnParseFormGloss{a}{present active subjunctive 1st plural}{ὁμολογέω}{to confess} τὰς ἁμαρτίας ἡμῶν, πιστός ἐστιν\FnParse{b}{present active indicative 3rd singular} καὶ δίκαιος ἵνα ἀφῇ\FnParse{c}{aorist active subjunctive 3rd singular} ἡμῖν τὰς ἁμαρτίας καὶ καθαρίσῃ\FnParse{d}{aorist active subjunctive 3rd singular} ἡμᾶς ἀπὸ πάσης ἀδικίας.\FnParseFormGloss{e}{genitive feminine singular}{ἀδικία}{wickedness} 
\MainTextVerseMark{1}{10}ἐὰν εἴπωμεν\FnParse{a}{aorist active subjunctive 1st plural} ὅτι οὐχ ἡμαρτήκαμεν,\FnParse{b}{perfect active indicative 1st plural} ψεύστην\FnParseFormGloss{c}{accusative masculine singular}{ψεύστης}{liar} ποιοῦμεν\FnParse{d}{present active indicative 1st plural} αὐτὸν καὶ ὁ λόγος αὐτοῦ οὐκ ἔστιν\FnParse{e}{present active indicative 3rd singular} ἐν ἡμῖν. 
\ChapterHeader{2}{Chapter 2}

\MainTextVerseMark{2}{1}Τεκνία\FnParseFormGloss{a}{neuter plural}{τεκνίον}{dear children} μου, ταῦτα γράφω\FnParse{b}{present active indicative 1st singular} ὑμῖν ἵνα μὴ ἁμάρτητε.\FnParse{c}{aorist active subjunctive 2nd plural} καὶ ἐάν τις ἁμάρτῃ,\FnParse{d}{aorist active subjunctive 3rd singular} παράκλητον\FnParseFormGloss{e}{accusative masculine singular}{παράκλητος}{counselor} ἔχομεν\FnParse{f}{present active indicative 1st plural} πρὸς τὸν πατέρα Ἰησοῦν Χριστὸν δίκαιον, 
\MainTextVerseMark{2}{2}καὶ αὐτὸς ἱλασμός\FnParseFormGloss{a}{nominative masculine singular}{ἱλασμός}{atoning sacrifice} ἐστιν\FnParse{b}{present active indicative 3rd singular} περὶ τῶν ἁμαρτιῶν ἡμῶν, οὐ περὶ τῶν ἡμετέρων\FnParseFormGloss{c}{genitive feminine plural}{ἡμέτερος}{our} δὲ μόνον ἀλλὰ καὶ περὶ ὅλου τοῦ κόσμου. 
\MainTextVerseMark{2}{3}Καὶ ἐν τούτῳ γινώσκομεν\FnParse{a}{present active indicative 1st plural} ὅτι ἐγνώκαμεν\FnParse{b}{perfect active indicative 1st plural} αὐτόν, ἐὰν τὰς ἐντολὰς αὐτοῦ τηρῶμεν.\FnParse{c}{present active subjunctive 1st plural} 
\MainTextVerseMark{2}{4}ὁ λέγων\FnParse{a}{present active participle nominative masculine singular} ⸀ὅτι Ἔγνωκα\FnParse{b}{perfect active indicative 1st singular} αὐτὸν καὶ τὰς ἐντολὰς αὐτοῦ μὴ τηρῶν\FnParse{c}{present active participle nominative masculine singular} ψεύστης\FnParseFormGloss{d}{nominative masculine singular}{ψεύστης}{liar} ἐστίν,\FnParse{e}{present active indicative 3rd singular} καὶ ἐν τούτῳ ἡ ἀλήθεια οὐκ ἔστιν·\FnParse{f}{present active indicative 3rd singular} 
\MainTextVerseMark{2}{5}ὃς δ’ ἂν τηρῇ\FnParse{a}{present active subjunctive 3rd singular} αὐτοῦ τὸν λόγον, ἀληθῶς\FnFormGloss{b}{ἀληθῶς}{truly} ἐν τούτῳ ἡ ἀγάπη τοῦ θεοῦ τετελείωται.\FnParseFormGloss{c}{perfect passive indicative 3rd singular}{τελειόω}{to perfect} ἐν τούτῳ γινώσκομεν\FnParse{d}{present active indicative 1st plural} ὅτι ἐν αὐτῷ ἐσμεν·\FnParse{e}{present active indicative 1st plural} 
\MainTextVerseMark{2}{6}ὁ λέγων\FnParse{a}{present active participle nominative masculine singular} ἐν αὐτῷ μένειν\FnParse{b}{present active infinitive} ὀφείλει\FnParse{c}{present active indicative 3rd singular} καθὼς ἐκεῖνος περιεπάτησεν\FnParse{d}{aorist active indicative 3rd singular} καὶ ⸀αὐτὸς περιπατεῖν.\FnParse{e}{present active infinitive} 
\MainTextVerseMark{2}{7}⸀Ἀγαπητοί, οὐκ ἐντολὴν καινὴν γράφω\FnParse{a}{present active indicative 1st singular} ὑμῖν, ἀλλ’ ἐντολὴν παλαιὰν\FnParseFormGloss{b}{accusative feminine singular}{παλαιός}{old} ἣν εἴχετε\FnParse{c}{imperfect active indicative 2nd plural} ἀπ’ ἀρχῆς· ἡ ἐντολὴ ἡ παλαιά\FnParseFormGloss{d}{nominative feminine singular}{παλαιός}{old} ἐστιν\FnParse{e}{present active indicative 3rd singular} ὁ λόγος ὃν ⸀ἠκούσατε.\FnParse{f}{aorist active indicative 2nd plural} 
\MainTextVerseMark{2}{8}πάλιν ἐντολὴν καινὴν γράφω\FnParse{a}{present active indicative 1st singular} ὑμῖν, ὅ ἐστιν\FnParse{b}{present active indicative 3rd singular} ἀληθὲς\FnParseFormGloss{c}{nominative neuter singular}{ἀληθής}{true} ἐν αὐτῷ καὶ ἐν ὑμῖν, ὅτι ἡ σκοτία\FnParseFormGloss{d}{nominative feminine singular}{σκοτία}{darkness} παράγεται\FnParseFormGloss{e}{present passive indicative 3rd singular}{παράγω}{to pass by} καὶ τὸ φῶς τὸ ἀληθινὸν\FnParseFormGloss{f}{nominative neuter singular}{ἀληθινός}{true} ἤδη φαίνει.\FnParse{g}{present active indicative 3rd singular} 
\MainTextVerseMark{2}{9}ὁ λέγων\FnParse{a}{present active participle nominative masculine singular} ἐν τῷ φωτὶ εἶναι\FnParse{b}{present active infinitive} καὶ τὸν ἀδελφὸν αὐτοῦ μισῶν\FnParse{c}{present active participle nominative masculine singular} ἐν τῇ σκοτίᾳ\FnParseFormGloss{d}{dative feminine singular}{σκοτία}{darkness} ἐστὶν\FnParse{e}{present active indicative 3rd singular} ἕως ἄρτι. 
\MainTextVerseMark{2}{10}ὁ ἀγαπῶν\FnParse{a}{present active participle nominative masculine singular} τὸν ἀδελφὸν αὐτοῦ ἐν τῷ φωτὶ μένει,\FnParse{b}{present active indicative 3rd singular} καὶ σκάνδαλον\FnParseFormGloss{c}{nominative neuter singular}{σκάνδαλον}{stumbling block} ἐν αὐτῷ οὐκ ἔστιν·\FnParse{d}{present active indicative 3rd singular} 
\MainTextVerseMark{2}{11}ὁ δὲ μισῶν\FnParse{a}{present active participle nominative masculine singular} τὸν ἀδελφὸν αὐτοῦ ἐν τῇ σκοτίᾳ\FnParseFormGloss{b}{dative feminine singular}{σκοτία}{darkness} ἐστὶν\FnParse{c}{present active indicative 3rd singular} καὶ ἐν τῇ σκοτίᾳ\FnParseFormGloss{d}{dative feminine singular}{σκοτία}{darkness} περιπατεῖ,\FnParse{e}{present active indicative 3rd singular} καὶ οὐκ οἶδεν\FnParse{f}{perfect active indicative 3rd singular} ποῦ ὑπάγει,\FnParse{g}{present active indicative 3rd singular} ὅτι ἡ σκοτία\FnParseFormGloss{h}{nominative feminine singular}{σκοτία}{darkness} ἐτύφλωσεν\FnParseFormGloss{i}{aorist active indicative 3rd singular}{τυφλόω}{to cause blindness} τοὺς ὀφθαλμοὺς αὐτοῦ. 
\MainTextVerseMark{2}{12}Γράφω\FnParse{a}{present active indicative 1st singular} ὑμῖν, τεκνία,\FnParseFormGloss{b}{neuter plural}{τεκνίον}{dear children} ὅτι ἀφέωνται\FnParse{c}{perfect passive indicative 3rd plural} ὑμῖν αἱ ἁμαρτίαι διὰ τὸ ὄνομα αὐτοῦ· 
\MainTextVerseMark{2}{13}γράφω\FnParse{a}{present active indicative 1st singular} ὑμῖν, πατέρες, ὅτι ἐγνώκατε\FnParse{b}{perfect active indicative 2nd plural} τὸν ἀπ’ ἀρχῆς· γράφω\FnParse{c}{present active indicative 1st singular} ὑμῖν, νεανίσκοι,\FnParseFormGloss{d}{masculine plural}{νεανίσκος}{young man} ὅτι νενικήκατε\FnParseFormGloss{e}{perfect active indicative 2nd plural}{νικάω}{to overcome} τὸν πονηρόν. 
\MainTextVerseMark{2}{14}⸀ἔγραψα\FnParse{a}{aorist active indicative 1st singular} ὑμῖν, παιδία, ὅτι ἐγνώκατε\FnParse{b}{perfect active indicative 2nd plural} τὸν πατέρα· ἔγραψα\FnParse{c}{aorist active indicative 1st singular} ὑμῖν, πατέρες, ὅτι ἐγνώκατε\FnParse{d}{perfect active indicative 2nd plural} τὸν ἀπ’ ἀρχῆς· ἔγραψα\FnParse{e}{aorist active indicative 1st singular} ὑμῖν, νεανίσκοι,\FnParseFormGloss{f}{masculine plural}{νεανίσκος}{young man} ὅτι ἰσχυροί\FnParseFormGloss{g}{nominative masculine plural}{ἰσχυρός}{powerful} ἐστε\FnParse{h}{present active indicative 2nd plural} καὶ ὁ λόγος τοῦ θεοῦ ἐν ὑμῖν μένει\FnParse{i}{present active indicative 3rd singular} καὶ νενικήκατε\FnParseFormGloss{j}{perfect active indicative 2nd plural}{νικάω}{to overcome} τὸν πονηρόν. 
\MainTextVerseMark{2}{15}Μὴ ἀγαπᾶτε\FnParse{a}{present active imperative 2nd plural} τὸν κόσμον μηδὲ τὰ ἐν τῷ κόσμῳ. ἐάν τις ἀγαπᾷ\FnParse{b}{present active subjunctive 3rd singular} τὸν κόσμον, οὐκ ἔστιν\FnParse{c}{present active indicative 3rd singular} ἡ ἀγάπη τοῦ πατρὸς ἐν αὐτῷ· 
\MainTextVerseMark{2}{16}ὅτι πᾶν τὸ ἐν τῷ κόσμῳ, ἡ ἐπιθυμία τῆς σαρκὸς καὶ ἡ ἐπιθυμία τῶν ὀφθαλμῶν καὶ ἡ ἀλαζονεία\FnParseFormGloss{a}{nominative feminine singular}{ἀλαζονεία}{boasting} τοῦ βίου,\FnParseFormGloss{b}{genitive masculine singular}{βίος}{life} οὐκ ἔστιν\FnParse{c}{present active indicative 3rd singular} ἐκ τοῦ πατρός, ἀλλὰ ἐκ τοῦ κόσμου ἐστίν·\FnParse{d}{present active indicative 3rd singular} 
\MainTextVerseMark{2}{17}καὶ ὁ κόσμος παράγεται\FnParseFormGloss{a}{present passive indicative 3rd singular}{παράγω}{to pass by} καὶ ἡ ἐπιθυμία αὐτοῦ, ὁ δὲ ποιῶν\FnParse{b}{present active participle nominative masculine singular} τὸ θέλημα τοῦ θεοῦ μένει\FnParse{c}{present active indicative 3rd singular} εἰς τὸν αἰῶνα. 
\MainTextVerseMark{2}{18}Παιδία, ἐσχάτη ὥρα ἐστίν,\FnParse{a}{present active indicative 3rd singular} καὶ καθὼς ἠκούσατε\FnParse{b}{aorist active indicative 2nd plural} ⸀ὅτι ἀντίχριστος\FnParseFormGloss{c}{nominative masculine singular}{ἀντίχριστος}{antichrist} ἔρχεται,\FnParse{d}{present middle indicative 3rd singular} καὶ νῦν ἀντίχριστοι\FnParseFormGloss{e}{nominative masculine plural}{ἀντίχριστος}{antichrist} πολλοὶ γεγόνασιν·\FnParse{f}{perfect active indicative 3rd plural} ὅθεν\FnFormGloss{g}{ὅθεν}{from where} γινώσκομεν\FnParse{h}{present active indicative 1st plural} ὅτι ἐσχάτη ὥρα ἐστίν.\FnParse{i}{present active indicative 3rd singular} 
\MainTextVerseMark{2}{19}ἐξ ἡμῶν ἐξῆλθαν,\FnParse{a}{aorist active indicative 3rd plural} ἀλλ’ οὐκ ἦσαν\FnParse{b}{imperfect active indicative 3rd plural} ἐξ ἡμῶν· εἰ γὰρ ⸂ἐξ ἡμῶν ἦσαν⸃,\FnParse{c}{imperfect active indicative 3rd plural} μεμενήκεισαν\FnParse{d}{pluperfect active indicative 3rd plural} ἂν μεθ’ ἡμῶν· ἀλλ’ ἵνα φανερωθῶσιν\FnParse{e}{aorist passive subjunctive 3rd plural} ὅτι οὐκ εἰσὶν\FnParse{f}{present active indicative 3rd plural} πάντες ἐξ ἡμῶν. 
\MainTextVerseMark{2}{20}καὶ ὑμεῖς χρῖσμα\FnParseFormGloss{a}{accusative neuter singular}{χρῖσμα}{anointing} ἔχετε\FnParse{b}{present active indicative 2nd plural} ἀπὸ τοῦ ἁγίου ⸀καὶ οἴδατε\FnParse{c}{perfect active indicative 2nd plural} ⸀πάντες· 
\MainTextVerseMark{2}{21}οὐκ ἔγραψα\FnParse{a}{aorist active indicative 1st singular} ὑμῖν ὅτι οὐκ οἴδατε\FnParse{b}{perfect active indicative 2nd plural} τὴν ἀλήθειαν, ἀλλ’ ὅτι οἴδατε\FnParse{c}{perfect active indicative 2nd plural} αὐτήν, καὶ ὅτι πᾶν ψεῦδος\FnParseFormGloss{d}{nominative neuter singular}{ψεῦδος}{lie} ἐκ τῆς ἀληθείας οὐκ ἔστιν.\FnParse{e}{present active indicative 3rd singular} 
\MainTextVerseMark{2}{22}τίς ἐστιν\FnParse{a}{present active indicative 3rd singular} ὁ ψεύστης\FnParseFormGloss{b}{nominative masculine singular}{ψεύστης}{liar} εἰ μὴ ὁ ἀρνούμενος\FnParse{c}{present middle participle nominative masculine singular} ὅτι Ἰησοῦς οὐκ ἔστιν\FnParse{d}{present active indicative 3rd singular} ὁ χριστός; οὗτός ἐστιν\FnParse{e}{present active indicative 3rd singular} ὁ ἀντίχριστος,\FnParseFormGloss{f}{nominative masculine singular}{ἀντίχριστος}{antichrist} ὁ ἀρνούμενος\FnParse{g}{present middle participle nominative masculine singular} τὸν πατέρα καὶ τὸν υἱόν. 
\MainTextVerseMark{2}{23}πᾶς ὁ ἀρνούμενος\FnParse{a}{present middle participle nominative masculine singular} τὸν υἱὸν οὐδὲ τὸν πατέρα ἔχει·\FnParse{b}{present active indicative 3rd singular} ⸂ὁ ὁμολογῶν\FnParseFormGloss{c}{present active participle nominative masculine singular}{ὁμολογέω}{to confess} τὸν υἱὸν καὶ τὸν πατέρα ἔχει⸃.\FnParse{d}{present active indicative 3rd singular} 
\MainTextVerseMark{2}{24}⸀ὑμεῖς ὃ ἠκούσατε\FnParse{a}{aorist active indicative 2nd plural} ἀπ’ ἀρχῆς, ἐν ὑμῖν μενέτω·\FnParse{b}{present active imperative 3rd singular} ἐὰν ἐν ὑμῖν μείνῃ\FnParse{c}{aorist active subjunctive 3rd singular} ὃ ἀπ’ ἀρχῆς ἠκούσατε,\FnParse{d}{aorist active indicative 2nd plural} καὶ ὑμεῖς ἐν τῷ υἱῷ καὶ ἐν τῷ πατρὶ μενεῖτε.\FnParse{e}{future active indicative 2nd plural} 
\MainTextVerseMark{2}{25}καὶ αὕτη ἐστὶν\FnParse{a}{present active indicative 3rd singular} ἡ ἐπαγγελία ἣν αὐτὸς ἐπηγγείλατο\FnParseFormGloss{b}{aorist middle indicative 3rd singular}{ἐπαγγέλλομαι}{to promise} ἡμῖν, τὴν ζωὴν τὴν αἰώνιον. 
\MainTextVerseMark{2}{26}Ταῦτα ἔγραψα\FnParse{a}{aorist active indicative 1st singular} ὑμῖν περὶ τῶν πλανώντων\FnParse{b}{present active participle genitive masculine plural} ὑμᾶς. 
\MainTextVerseMark{2}{27}καὶ ὑμεῖς τὸ χρῖσμα\FnParseFormGloss{a}{nominative neuter singular}{χρῖσμα}{anointing} ὃ ἐλάβετε\FnParse{b}{aorist active indicative 2nd plural} ἀπ’ αὐτοῦ ⸂μένει\FnParse{c}{present active indicative 3rd singular} ἐν ὑμῖν⸃, καὶ οὐ χρείαν ἔχετε\FnParse{d}{present active indicative 2nd plural} ἵνα τις διδάσκῃ\FnParse{e}{present active subjunctive 3rd singular} ὑμᾶς· ἀλλ’ ὡς τὸ ⸀αὐτοῦ χρῖσμα\FnParseFormGloss{f}{nominative neuter singular}{χρῖσμα}{anointing} διδάσκει\FnParse{g}{present active indicative 3rd singular} ὑμᾶς περὶ πάντων, καὶ ἀληθές\FnParseFormGloss{h}{nominative neuter singular}{ἀληθής}{true} ἐστιν\FnParse{i}{present active indicative 3rd singular} καὶ οὐκ ἔστιν\FnParse{j}{present active indicative 3rd singular} ψεῦδος,\FnParseFormGloss{k}{nominative neuter singular}{ψεῦδος}{lie} καὶ καθὼς ἐδίδαξεν\FnParse{l}{aorist active indicative 3rd singular} ὑμᾶς, ⸀μένετε\FnParse{m}{present active indicative 2nd plural} ἐν αὐτῷ. 
\MainTextVerseMark{2}{28}Καὶ νῦν, τεκνία,\FnParseFormGloss{a}{neuter plural}{τεκνίον}{dear children} μένετε\FnParse{b}{present active imperative 2nd plural} ἐν αὐτῷ, ἵνα ⸀ἐὰν φανερωθῇ\FnParse{c}{aorist passive subjunctive 3rd singular} ⸀σχῶμεν\FnParse{d}{aorist active subjunctive 1st plural} παρρησίαν καὶ μὴ αἰσχυνθῶμεν\FnParseFormGloss{e}{aorist passive subjunctive 1st plural}{αἰσχύνομαι}{to be ashamed} ἀπ’ αὐτοῦ ἐν τῇ παρουσίᾳ\FnParseFormGloss{f}{dative feminine singular}{παρουσία}{presence} αὐτοῦ. 
\MainTextVerseMark{2}{29}ἐὰν εἰδῆτε\FnParse{a}{perfect active subjunctive 2nd plural} ὅτι δίκαιός ἐστιν,\FnParse{b}{present active indicative 3rd singular} γινώσκετε\FnParse{c}{present active indicative 2nd plural} ⸀ὅτι πᾶς ὁ ποιῶν\FnParse{d}{present active participle nominative masculine singular} τὴν δικαιοσύνην ἐξ αὐτοῦ γεγέννηται.\FnParse{e}{perfect passive indicative 3rd singular} 
\ChapterHeader{3}{Chapter 3}

\MainTextVerseMark{3}{1}ἴδετε\FnParse{a}{aorist active imperative 2nd plural} ποταπὴν\FnParseFormGloss{b}{accusative feminine singular}{ποταπός}{of what kind?} ἀγάπην δέδωκεν\FnParse{c}{perfect active indicative 3rd singular} ἡμῖν ὁ πατὴρ ἵνα τέκνα θεοῦ κληθῶμεν,\FnParse{d}{aorist passive subjunctive 1st plural} ⸂καὶ ἐσμέν⸃.\FnParse{e}{present active indicative 1st plural} διὰ τοῦτο ὁ κόσμος οὐ γινώσκει\FnParse{f}{present active indicative 3rd singular} ⸀ἡμᾶς ὅτι οὐκ ἔγνω\FnParse{g}{aorist active indicative 3rd singular} αὐτόν. 
\MainTextVerseMark{3}{2}ἀγαπητοί, νῦν τέκνα θεοῦ ἐσμεν,\FnParse{a}{present active indicative 1st plural} καὶ οὔπω\FnFormGloss{b}{οὔπω}{not yet} ἐφανερώθη\FnParse{c}{aorist passive indicative 3rd singular} τί ἐσόμεθα.\FnParse{d}{future middle indicative 1st plural} ⸀οἴδαμεν\FnParse{e}{perfect active indicative 1st plural} ὅτι ἐὰν φανερωθῇ\FnParse{f}{aorist passive subjunctive 3rd singular} ὅμοιοι αὐτῷ ἐσόμεθα,\FnParse{g}{future middle indicative 1st plural} ὅτι ὀψόμεθα\FnParse{h}{future middle indicative 1st plural} αὐτὸν καθώς ἐστιν.\FnParse{i}{present active indicative 3rd singular} 
\MainTextVerseMark{3}{3}καὶ πᾶς ὁ ἔχων\FnParse{a}{present active participle nominative masculine singular} τὴν ἐλπίδα ταύτην ἐπ’ αὐτῷ ἁγνίζει\FnParseFormGloss{b}{present active indicative 3rd singular}{ἁγνίζω}{to purify} ἑαυτὸν καθὼς ἐκεῖνος ἁγνός\FnParseFormGloss{c}{nominative masculine singular}{ἁγνός}{pure} ἐστιν.\FnParse{d}{present active indicative 3rd singular} 
\MainTextVerseMark{3}{4}Πᾶς ὁ ποιῶν\FnParse{a}{present active participle nominative masculine singular} τὴν ἁμαρτίαν καὶ τὴν ἀνομίαν\FnParseFormGloss{b}{accusative feminine singular}{ἀνομία}{wickedness} ποιεῖ,\FnParse{c}{present active indicative 3rd singular} καὶ ἡ ἁμαρτία ἐστὶν\FnParse{d}{present active indicative 3rd singular} ἡ ἀνομία.\FnParseFormGloss{e}{nominative feminine singular}{ἀνομία}{wickedness} 
\MainTextVerseMark{3}{5}καὶ οἴδατε\FnParse{a}{perfect active indicative 2nd plural} ὅτι ἐκεῖνος ἐφανερώθη\FnParse{b}{aorist passive indicative 3rd singular} ἵνα τὰς ⸀ἁμαρτίας ἄρῃ,\FnParse{c}{aorist active subjunctive 3rd singular} καὶ ἁμαρτία ἐν αὐτῷ οὐκ ἔστιν.\FnParse{d}{present active indicative 3rd singular} 
\MainTextVerseMark{3}{6}πᾶς ὁ ἐν αὐτῷ μένων\FnParse{a}{present active participle nominative masculine singular} οὐχ ἁμαρτάνει·\FnParse{b}{present active indicative 3rd singular} πᾶς ὁ ἁμαρτάνων\FnParse{c}{present active participle nominative masculine singular} οὐχ ἑώρακεν\FnParse{d}{perfect active indicative 3rd singular} αὐτὸν οὐδὲ ἔγνωκεν\FnParse{e}{perfect active indicative 3rd singular} αὐτόν. 
\MainTextVerseMark{3}{7}τεκνία,\FnParseFormGloss{a}{neuter plural}{τεκνίον}{dear children} μηδεὶς πλανάτω\FnParse{b}{present active imperative 3rd singular} ὑμᾶς· ὁ ποιῶν\FnParse{c}{present active participle nominative masculine singular} τὴν δικαιοσύνην δίκαιός ἐστιν,\FnParse{d}{present active indicative 3rd singular} καθὼς ἐκεῖνος δίκαιός ἐστιν·\FnParse{e}{present active indicative 3rd singular} 
\MainTextVerseMark{3}{8}ὁ ποιῶν\FnParse{a}{present active participle nominative masculine singular} τὴν ἁμαρτίαν ἐκ τοῦ διαβόλου ἐστίν,\FnParse{b}{present active indicative 3rd singular} ὅτι ἀπ’ ἀρχῆς ὁ διάβολος ἁμαρτάνει.\FnParse{c}{present active indicative 3rd singular} εἰς τοῦτο ἐφανερώθη\FnParse{d}{aorist passive indicative 3rd singular} ὁ υἱὸς τοῦ θεοῦ ἵνα λύσῃ\FnParse{e}{aorist active subjunctive 3rd singular} τὰ ἔργα τοῦ διαβόλου. 
\MainTextVerseMark{3}{9}πᾶς ὁ γεγεννημένος\FnParse{a}{perfect passive participle nominative masculine singular} ἐκ τοῦ θεοῦ ἁμαρτίαν οὐ ποιεῖ,\FnParse{b}{present active indicative 3rd singular} ὅτι σπέρμα αὐτοῦ ἐν αὐτῷ μένει,\FnParse{c}{present active indicative 3rd singular} καὶ οὐ δύναται\FnParse{d}{present middle indicative 3rd singular} ἁμαρτάνειν,\FnParse{e}{present active infinitive} ὅτι ἐκ τοῦ θεοῦ γεγέννηται.\FnParse{f}{perfect passive indicative 3rd singular} 
\MainTextVerseMark{3}{10}ἐν τούτῳ φανερά\FnParseFormGloss{a}{nominative neuter plural}{φανερός}{visible} ἐστιν\FnParse{b}{present active indicative 3rd singular} τὰ τέκνα τοῦ θεοῦ καὶ τὰ τέκνα τοῦ διαβόλου· πᾶς ὁ μὴ ποιῶν\FnParse{c}{present active participle nominative masculine singular} δικαιοσύνην οὐκ ἔστιν\FnParse{d}{present active indicative 3rd singular} ἐκ τοῦ θεοῦ, καὶ ὁ μὴ ἀγαπῶν\FnParse{e}{present active participle nominative masculine singular} τὸν ἀδελφὸν αὐτοῦ. 
\MainTextVerseMark{3}{11}Ὅτι αὕτη ἐστὶν\FnParse{a}{present active indicative 3rd singular} ἡ ἀγγελία\FnParseFormGloss{b}{nominative feminine singular}{ἀγγελία}{message} ἣν ἠκούσατε\FnParse{c}{aorist active indicative 2nd plural} ἀπ’ ἀρχῆς, ἵνα ἀγαπῶμεν\FnParse{d}{present active subjunctive 1st plural} ἀλλήλους· 
\MainTextVerseMark{3}{12}οὐ καθὼς Κάϊν\FnParseFormGloss{a}{nominative masculine singular}{Κάϊν}{Cain} ἐκ τοῦ πονηροῦ ἦν\FnParse{b}{imperfect active indicative 3rd singular} καὶ ἔσφαξεν\FnParseFormGloss{c}{aorist active indicative 3rd singular}{σφάζω}{to kill} τὸν ἀδελφὸν αὐτοῦ· καὶ χάριν\FnFormGloss{d}{χάριν}{therefore} τίνος ἔσφαξεν\FnParseFormGloss{e}{aorist active indicative 3rd singular}{σφάζω}{to kill} αὐτόν; ὅτι τὰ ἔργα αὐτοῦ πονηρὰ ἦν,\FnParse{f}{imperfect active indicative 3rd singular} τὰ δὲ τοῦ ἀδελφοῦ αὐτοῦ δίκαια. 
\MainTextVerseMark{3}{13}⸀μὴ θαυμάζετε,\FnParse{a}{present active imperative 2nd plural} ⸀ἀδελφοί, εἰ μισεῖ\FnParse{b}{present active indicative 3rd singular} ὑμᾶς ὁ κόσμος. 
\MainTextVerseMark{3}{14}ἡμεῖς οἴδαμεν\FnParse{a}{perfect active indicative 1st plural} ὅτι μεταβεβήκαμεν\FnParseFormGloss{b}{perfect active indicative 1st plural}{μεταβαίνω}{to go on} ἐκ τοῦ θανάτου εἰς τὴν ζωήν, ὅτι ἀγαπῶμεν\FnParse{c}{present active indicative 1st plural} τοὺς ἀδελφούς· ὁ μὴ ⸀ἀγαπῶν\FnParse{d}{present active participle nominative masculine singular} μένει\FnParse{e}{present active indicative 3rd singular} ἐν τῷ θανάτῳ. 
\MainTextVerseMark{3}{15}πᾶς ὁ μισῶν\FnParse{a}{present active participle nominative masculine singular} τὸν ἀδελφὸν αὐτοῦ ἀνθρωποκτόνος\FnParseFormGloss{b}{nominative masculine singular}{ἀνθρωποκτόνος}{murderer} ἐστίν,\FnParse{c}{present active indicative 3rd singular} καὶ οἴδατε\FnParse{d}{perfect active indicative 2nd plural} ὅτι πᾶς ἀνθρωποκτόνος\FnParseFormGloss{e}{nominative masculine singular}{ἀνθρωποκτόνος}{murderer} οὐκ ἔχει\FnParse{f}{present active indicative 3rd singular} ζωὴν αἰώνιον ἐν ⸀αὐτῷ μένουσαν.\FnParse{g}{present active participle accusative feminine singular} 
\MainTextVerseMark{3}{16}ἐν τούτῳ ἐγνώκαμεν\FnParse{a}{perfect active indicative 1st plural} τὴν ἀγάπην, ὅτι ἐκεῖνος ὑπὲρ ἡμῶν τὴν ψυχὴν αὐτοῦ ἔθηκεν·\FnParse{b}{aorist active indicative 3rd singular} καὶ ἡμεῖς ὀφείλομεν\FnParse{c}{present active indicative 1st plural} ὑπὲρ τῶν ἀδελφῶν τὰς ψυχὰς ⸀θεῖναι.\FnParse{d}{aorist active infinitive} 
\MainTextVerseMark{3}{17}ὃς δ’ ἂν ἔχῃ\FnParse{a}{present active subjunctive 3rd singular} τὸν βίον\FnParseFormGloss{b}{accusative masculine singular}{βίος}{life} τοῦ κόσμου καὶ θεωρῇ\FnParse{c}{present active subjunctive 3rd singular} τὸν ἀδελφὸν αὐτοῦ χρείαν ἔχοντα\FnParse{d}{present active participle accusative masculine singular} καὶ κλείσῃ\FnParseFormGloss{e}{aorist active subjunctive 3rd singular}{κλείω}{to close} τὰ σπλάγχνα\FnParseFormGloss{f}{accusative neuter plural}{σπλάγχνον}{intestines} αὐτοῦ ἀπ’ αὐτοῦ, πῶς ἡ ἀγάπη τοῦ θεοῦ μένει\FnParse{g}{present active indicative 3rd singular} ἐν αὐτῷ; 
\MainTextVerseMark{3}{18}⸀Τεκνία,\FnParseFormGloss{a}{neuter plural}{τεκνίον}{dear children} μὴ ἀγαπῶμεν\FnParse{b}{present active subjunctive 1st plural} λόγῳ μηδὲ τῇ γλώσσῃ ἀλλὰ ἐν ἔργῳ καὶ ἀληθείᾳ. 
\MainTextVerseMark{3}{19}⸀ἐν τούτῳ ⸀γνωσόμεθα\FnParse{a}{future middle indicative 1st plural} ὅτι ἐκ τῆς ἀληθείας ἐσμέν,\FnParse{b}{present active indicative 1st plural} καὶ ἔμπροσθεν αὐτοῦ πείσομεν\FnParse{c}{future active indicative 1st plural} ⸂τὴν καρδίαν⸃ ἡμῶν 
\MainTextVerseMark{3}{20}ὅτι ἐὰν καταγινώσκῃ\FnParseFormGloss{a}{present active subjunctive 3rd singular}{καταγινώσκω}{to condemn} ἡμῶν ἡ καρδία, ὅτι μείζων ἐστὶν\FnParse{b}{present active indicative 3rd singular} ὁ θεὸς τῆς καρδίας ἡμῶν καὶ γινώσκει\FnParse{c}{present active indicative 3rd singular} πάντα. 
\MainTextVerseMark{3}{21}ἀγαπητοί, ἐὰν ἡ καρδία ⸂μὴ καταγινώσκῃ\FnParseFormGloss{a}{present active subjunctive 3rd singular}{καταγινώσκω}{to condemn} ἡμῶν⸃, παρρησίαν ἔχομεν\FnParse{b}{present active indicative 1st plural} πρὸς τὸν θεόν, 
\MainTextVerseMark{3}{22}καὶ ὃ ⸀ἐὰν αἰτῶμεν\FnParse{a}{present active subjunctive 1st plural} λαμβάνομεν\FnParse{b}{present active indicative 1st plural} ⸀ἀπ’ αὐτοῦ, ὅτι τὰς ἐντολὰς αὐτοῦ τηροῦμεν\FnParse{c}{present active indicative 1st plural} καὶ τὰ ἀρεστὰ\FnParseFormGloss{d}{accusative neuter plural}{ἀρεστός}{pleasing} ἐνώπιον αὐτοῦ ποιοῦμεν.\FnParse{e}{present active indicative 1st plural} 
\MainTextVerseMark{3}{23}καὶ αὕτη ἐστὶν\FnParse{a}{present active indicative 3rd singular} ἡ ἐντολὴ αὐτοῦ, ἵνα ⸀πιστεύσωμεν\FnParse{b}{aorist active subjunctive 1st plural} τῷ ὀνόματι τοῦ υἱοῦ αὐτοῦ Ἰησοῦ Χριστοῦ καὶ ἀγαπῶμεν\FnParse{c}{present active subjunctive 1st plural} ἀλλήλους, καθὼς ἔδωκεν\FnParse{d}{aorist active indicative 3rd singular} ἐντολὴν ⸀ἡμῖν. 
\MainTextVerseMark{3}{24}καὶ ὁ τηρῶν\FnParse{a}{present active participle nominative masculine singular} τὰς ἐντολὰς αὐτοῦ ἐν αὐτῷ μένει\FnParse{b}{present active indicative 3rd singular} καὶ αὐτὸς ἐν αὐτῷ· καὶ ἐν τούτῳ γινώσκομεν\FnParse{c}{present active indicative 1st plural} ὅτι μένει\FnParse{d}{present active indicative 3rd singular} ἐν ἡμῖν, ἐκ τοῦ πνεύματος οὗ ἡμῖν ἔδωκεν.\FnParse{e}{aorist active indicative 3rd singular} 
\ChapterHeader{4}{Chapter 4}

\MainTextVerseMark{4}{1}Ἀγαπητοί, μὴ παντὶ πνεύματι πιστεύετε,\FnParse{a}{present active imperative 2nd plural} ἀλλὰ δοκιμάζετε\FnParseFormGloss{b}{present active imperative 2nd plural}{δοκιμάζω}{to test} τὰ πνεύματα εἰ ἐκ τοῦ θεοῦ ἐστιν,\FnParse{c}{present active indicative 3rd singular} ὅτι πολλοὶ ψευδοπροφῆται\FnParseFormGloss{d}{nominative masculine plural}{ψευδοπροφήτης}{false prophet} ἐξεληλύθασιν\FnParse{e}{perfect active indicative 3rd plural} εἰς τὸν κόσμον. 
\MainTextVerseMark{4}{2}ἐν τούτῳ ⸀γινώσκετε\FnParse{a}{present active indicative 2nd plural} τὸ πνεῦμα τοῦ θεοῦ· πᾶν πνεῦμα ὃ ὁμολογεῖ\FnParseFormGloss{b}{present active indicative 3rd singular}{ὁμολογέω}{to confess} Ἰησοῦν Χριστὸν ἐν σαρκὶ ἐληλυθότα\FnParse{c}{perfect active participle accusative masculine singular} ἐκ τοῦ θεοῦ ἐστιν,\FnParse{d}{present active indicative 3rd singular} 
\MainTextVerseMark{4}{3}καὶ πᾶν πνεῦμα ὃ μὴ ὁμολογεῖ\FnParseFormGloss{a}{present active indicative 3rd singular}{ὁμολογέω}{to confess} ⸀τὸν ⸀Ἰησοῦν ἐκ τοῦ θεοῦ οὐκ ἔστιν·\FnParse{b}{present active indicative 3rd singular} καὶ τοῦτό ἐστιν\FnParse{c}{present active indicative 3rd singular} τὸ τοῦ ἀντιχρίστου,\FnParseFormGloss{d}{genitive masculine singular}{ἀντίχριστος}{antichrist} ὃ ἀκηκόατε\FnParse{e}{perfect active indicative 2nd plural} ὅτι ἔρχεται,\FnParse{f}{present middle indicative 3rd singular} καὶ νῦν ἐν τῷ κόσμῳ ἐστὶν\FnParse{g}{present active indicative 3rd singular} ἤδη. 
\MainTextVerseMark{4}{4}ὑμεῖς ἐκ τοῦ θεοῦ ἐστε,\FnParse{a}{present active indicative 2nd plural} τεκνία,\FnParseFormGloss{b}{neuter plural}{τεκνίον}{dear children} καὶ νενικήκατε\FnParseFormGloss{c}{perfect active indicative 2nd plural}{νικάω}{to overcome} αὐτούς, ὅτι μείζων ἐστὶν\FnParse{d}{present active indicative 3rd singular} ὁ ἐν ὑμῖν ἢ ὁ ἐν τῷ κόσμῳ· 
\MainTextVerseMark{4}{5}αὐτοὶ ἐκ τοῦ κόσμου εἰσίν·\FnParse{a}{present active indicative 3rd plural} διὰ τοῦτο ἐκ τοῦ κόσμου λαλοῦσιν\FnParse{b}{present active indicative 3rd plural} καὶ ὁ κόσμος αὐτῶν ἀκούει.\FnParse{c}{present active indicative 3rd singular} 
\MainTextVerseMark{4}{6}ἡμεῖς ἐκ τοῦ θεοῦ ἐσμεν·\FnParse{a}{present active indicative 1st plural} ὁ γινώσκων\FnParse{b}{present active participle nominative masculine singular} τὸν θεὸν ἀκούει\FnParse{c}{present active indicative 3rd singular} ἡμῶν, ὃς οὐκ ἔστιν\FnParse{d}{present active indicative 3rd singular} ἐκ τοῦ θεοῦ οὐκ ἀκούει\FnParse{e}{present active indicative 3rd singular} ἡμῶν. ἐκ τούτου γινώσκομεν\FnParse{f}{present active indicative 1st plural} τὸ πνεῦμα τῆς ἀληθείας καὶ τὸ πνεῦμα τῆς πλάνης.\FnParseFormGloss{g}{genitive feminine singular}{πλάνη}{error} 
\MainTextVerseMark{4}{7}Ἀγαπητοί, ἀγαπῶμεν\FnParse{a}{present active subjunctive 1st plural} ἀλλήλους, ὅτι ἡ ἀγάπη ἐκ τοῦ θεοῦ ἐστιν,\FnParse{b}{present active indicative 3rd singular} καὶ πᾶς ὁ ἀγαπῶν\FnParse{c}{present active participle nominative masculine singular} ἐκ τοῦ θεοῦ γεγέννηται\FnParse{d}{perfect passive indicative 3rd singular} καὶ γινώσκει\FnParse{e}{present active indicative 3rd singular} τὸν θεόν. 
\MainTextVerseMark{4}{8}ὁ μὴ ἀγαπῶν\FnParse{a}{present active participle nominative masculine singular} οὐκ ἔγνω\FnParse{b}{aorist active indicative 3rd singular} τὸν θεόν, ὅτι ὁ θεὸς ἀγάπη ἐστίν.\FnParse{c}{present active indicative 3rd singular} 
\MainTextVerseMark{4}{9}ἐν τούτῳ ἐφανερώθη\FnParse{a}{aorist passive indicative 3rd singular} ἡ ἀγάπη τοῦ θεοῦ ἐν ἡμῖν, ὅτι τὸν υἱὸν αὐτοῦ τὸν μονογενῆ\FnParseFormGloss{b}{accusative masculine singular}{μονογενής}{one and only} ἀπέσταλκεν\FnParse{c}{perfect active indicative 3rd singular} ὁ θεὸς εἰς τὸν κόσμον ἵνα ζήσωμεν\FnParse{d}{aorist active subjunctive 1st plural} δι’ αὐτοῦ. 
\MainTextVerseMark{4}{10}ἐν τούτῳ ἐστὶν\FnParse{a}{present active indicative 3rd singular} ἡ ἀγάπη, οὐχ ὅτι ἡμεῖς ⸀ἠγαπήκαμεν\FnParse{b}{perfect active indicative 1st plural} τὸν θεόν, ἀλλ’ ὅτι αὐτὸς ἠγάπησεν\FnParse{c}{aorist active indicative 3rd singular} ἡμᾶς καὶ ἀπέστειλεν\FnParse{d}{aorist active indicative 3rd singular} τὸν υἱὸν αὐτοῦ ἱλασμὸν\FnParseFormGloss{e}{accusative masculine singular}{ἱλασμός}{atoning sacrifice} περὶ τῶν ἁμαρτιῶν ἡμῶν. 
\MainTextVerseMark{4}{11}ἀγαπητοί, εἰ οὕτως ὁ θεὸς ἠγάπησεν\FnParse{a}{aorist active indicative 3rd singular} ἡμᾶς, καὶ ἡμεῖς ὀφείλομεν\FnParse{b}{present active indicative 1st plural} ἀλλήλους ἀγαπᾶν.\FnParse{c}{present active infinitive} 
\MainTextVerseMark{4}{12}θεὸν οὐδεὶς πώποτε\FnFormGloss{a}{πώποτε}{ever} τεθέαται·\FnParseFormGloss{b}{perfect middle indicative 3rd singular}{θεάομαι}{to see} ἐὰν ἀγαπῶμεν\FnParse{c}{present active subjunctive 1st plural} ἀλλήλους, ὁ θεὸς ἐν ἡμῖν μένει\FnParse{d}{present active indicative 3rd singular} καὶ ἡ ἀγάπη αὐτοῦ ⸂ἐν ἡμῖν τετελειωμένη\FnParseFormGloss{e}{perfect passive participle nominative feminine singular}{τελειόω}{to perfect} ἐστιν⸃.\FnParse{f}{present active indicative 3rd singular} 
\MainTextVerseMark{4}{13}Ἐν τούτῳ γινώσκομεν\FnParse{a}{present active indicative 1st plural} ὅτι ἐν αὐτῷ μένομεν\FnParse{b}{present active indicative 1st plural} καὶ αὐτὸς ἐν ἡμῖν, ὅτι ἐκ τοῦ πνεύματος αὐτοῦ δέδωκεν\FnParse{c}{perfect active indicative 3rd singular} ἡμῖν. 
\MainTextVerseMark{4}{14}καὶ ἡμεῖς τεθεάμεθα\FnParseFormGloss{a}{perfect middle indicative 1st plural}{θεάομαι}{to see} καὶ μαρτυροῦμεν\FnParse{b}{present active indicative 1st plural} ὅτι ὁ πατὴρ ἀπέσταλκεν\FnParse{c}{perfect active indicative 3rd singular} τὸν υἱὸν σωτῆρα\FnParseFormGloss{d}{accusative masculine singular}{σωτήρ}{Savior} τοῦ κόσμου. 
\MainTextVerseMark{4}{15}ὃς ⸀ἐὰν ὁμολογήσῃ\FnParseFormGloss{a}{aorist active subjunctive 3rd singular}{ὁμολογέω}{to confess} ὅτι ⸀Ἰησοῦς ἐστιν\FnParse{b}{present active indicative 3rd singular} ὁ υἱὸς τοῦ θεοῦ, ὁ θεὸς ἐν αὐτῷ μένει\FnParse{c}{present active indicative 3rd singular} καὶ αὐτὸς ἐν τῷ θεῷ. 
\MainTextVerseMark{4}{16}καὶ ἡμεῖς ἐγνώκαμεν\FnParse{a}{perfect active indicative 1st plural} καὶ πεπιστεύκαμεν\FnParse{b}{perfect active indicative 1st plural} τὴν ἀγάπην ἣν ἔχει\FnParse{c}{present active indicative 3rd singular} ὁ θεὸς ἐν ἡμῖν. Ὁ θεὸς ἀγάπη ἐστίν,\FnParse{d}{present active indicative 3rd singular} καὶ ὁ μένων\FnParse{e}{present active participle nominative masculine singular} ἐν τῇ ἀγάπῃ ἐν τῷ θεῷ μένει\FnParse{f}{present active indicative 3rd singular} καὶ ὁ θεὸς ἐν αὐτῷ ⸀μένει.\FnParse{g}{present active indicative 3rd singular} 
\MainTextVerseMark{4}{17}ἐν τούτῳ τετελείωται\FnParseFormGloss{a}{perfect passive indicative 3rd singular}{τελειόω}{to perfect} ἡ ἀγάπη μεθ’ ἡμῶν, ἵνα παρρησίαν ἔχωμεν\FnParse{b}{present active subjunctive 1st plural} ἐν τῇ ἡμέρᾳ τῆς κρίσεως, ὅτι καθὼς ἐκεῖνός ἐστιν\FnParse{c}{present active indicative 3rd singular} καὶ ἡμεῖς ἐσμεν\FnParse{d}{present active indicative 1st plural} ἐν τῷ κόσμῳ τούτῳ. 
\MainTextVerseMark{4}{18}φόβος οὐκ ἔστιν\FnParse{a}{present active indicative 3rd singular} ἐν τῇ ἀγάπῃ, ἀλλ’ ἡ τελεία\FnParseFormGloss{b}{nominative feminine singular}{τέλειος}{perfect} ἀγάπη ἔξω βάλλει\FnParse{c}{present active indicative 3rd singular} τὸν φόβον, ὅτι ὁ φόβος κόλασιν\FnParseFormGloss{d}{accusative feminine singular}{κόλασις}{punishment} ἔχει,\FnParse{e}{present active indicative 3rd singular} ὁ δὲ φοβούμενος\FnParse{f}{present middle participle nominative masculine singular} οὐ τετελείωται\FnParseFormGloss{g}{perfect passive indicative 3rd singular}{τελειόω}{to perfect} ἐν τῇ ἀγάπῃ. 
\MainTextVerseMark{4}{19}ἡμεῖς ⸀ἀγαπῶμεν,\FnParse{a}{present active indicative 1st plural} ὅτι αὐτὸς πρῶτος ἠγάπησεν\FnParse{b}{aorist active indicative 3rd singular} ἡμᾶς. 
\MainTextVerseMark{4}{20}ἐάν τις εἴπῃ\FnParse{a}{aorist active subjunctive 3rd singular} ὅτι Ἀγαπῶ\FnParse{b}{present active indicative 1st singular} τὸν θεόν, καὶ τὸν ἀδελφὸν αὐτοῦ μισῇ,\FnParse{c}{present active subjunctive 3rd singular} ψεύστης\FnParseFormGloss{d}{nominative masculine singular}{ψεύστης}{liar} ἐστίν·\FnParse{e}{present active indicative 3rd singular} ὁ γὰρ μὴ ἀγαπῶν\FnParse{f}{present active participle nominative masculine singular} τὸν ἀδελφὸν αὐτοῦ ὃν ἑώρακεν,\FnParse{g}{perfect active indicative 3rd singular} τὸν θεὸν ὃν οὐχ ἑώρακεν\FnParse{h}{perfect active indicative 3rd singular} ⸀οὐ δύναται\FnParse{i}{present middle indicative 3rd singular} ἀγαπᾶν.\FnParse{j}{present active infinitive} 
\MainTextVerseMark{4}{21}καὶ ταύτην τὴν ἐντολὴν ἔχομεν\FnParse{a}{present active indicative 1st plural} ἀπ’ αὐτοῦ, ἵνα ὁ ἀγαπῶν\FnParse{b}{present active participle nominative masculine singular} τὸν θεὸν ἀγαπᾷ\FnParse{c}{present active subjunctive 3rd singular} καὶ τὸν ἀδελφὸν αὐτοῦ. 
\ChapterHeader{5}{Chapter 5}

\MainTextVerseMark{5}{1}Πᾶς ὁ πιστεύων\FnParse{a}{present active participle nominative masculine singular} ὅτι Ἰησοῦς ἐστιν\FnParse{b}{present active indicative 3rd singular} ὁ χριστὸς ἐκ τοῦ θεοῦ γεγέννηται,\FnParse{c}{perfect passive indicative 3rd singular} καὶ πᾶς ὁ ἀγαπῶν\FnParse{d}{present active participle nominative masculine singular} τὸν γεννήσαντα\FnParse{e}{aorist active participle accusative masculine singular} ἀγαπᾷ\FnParse{f}{present active indicative 3rd singular} ⸀καὶ τὸν γεγεννημένον\FnParse{g}{perfect passive participle accusative masculine singular} ἐξ αὐτοῦ. 
\MainTextVerseMark{5}{2}ἐν τούτῳ γινώσκομεν\FnParse{a}{present active indicative 1st plural} ὅτι ἀγαπῶμεν\FnParse{b}{present active indicative 1st plural} τὰ τέκνα τοῦ θεοῦ, ὅταν τὸν θεὸν ἀγαπῶμεν\FnParse{c}{present active subjunctive 1st plural} καὶ τὰς ἐντολὰς αὐτοῦ ⸀ποιῶμεν·\FnParse{d}{present active subjunctive 1st plural} 
\MainTextVerseMark{5}{3}αὕτη γάρ ἐστιν\FnParse{a}{present active indicative 3rd singular} ἡ ἀγάπη τοῦ θεοῦ ἵνα τὰς ἐντολὰς αὐτοῦ τηρῶμεν,\FnParse{b}{present active subjunctive 1st plural} καὶ αἱ ἐντολαὶ αὐτοῦ βαρεῖαι\FnParseFormGloss{c}{nominative feminine plural}{βαρύς}{burdensome} οὐκ εἰσίν,\FnParse{d}{present active indicative 3rd plural} 
\MainTextVerseMark{5}{4}ὅτι πᾶν τὸ γεγεννημένον\FnParse{a}{perfect passive participle nominative neuter singular} ἐκ τοῦ θεοῦ νικᾷ\FnParseFormGloss{b}{present active indicative 3rd singular}{νικάω}{to overcome} τὸν κόσμον. καὶ αὕτη ἐστὶν\FnParse{c}{present active indicative 3rd singular} ἡ νίκη\FnParseFormGloss{d}{nominative feminine singular}{νίκη}{victory} ἡ νικήσασα\FnParseFormGloss{e}{aorist active participle nominative feminine singular}{νικάω}{to overcome} τὸν κόσμον, ἡ πίστις ἡμῶν· 
\MainTextVerseMark{5}{5}τίς ⸂δέ ἐστιν⸃\FnParse{a}{present active indicative 3rd singular} ὁ νικῶν\FnParseFormGloss{b}{present active participle nominative masculine singular}{νικάω}{to overcome} τὸν κόσμον εἰ μὴ ὁ πιστεύων\FnParse{c}{present active participle nominative masculine singular} ὅτι Ἰησοῦς ἐστιν\FnParse{d}{present active indicative 3rd singular} ὁ υἱὸς τοῦ θεοῦ; 
\MainTextVerseMark{5}{6}Οὗτός ἐστιν\FnParse{a}{present active indicative 3rd singular} ὁ ἐλθὼν\FnParse{b}{aorist active participle nominative masculine singular} δι’ ὕδατος καὶ αἵματος, Ἰησοῦς Χριστός· οὐκ ἐν τῷ ὕδατι μόνον ἀλλ’ ἐν τῷ ὕδατι καὶ ⸀ἐν τῷ αἵματι· καὶ τὸ πνεῦμά ἐστιν\FnParse{c}{present active indicative 3rd singular} τὸ μαρτυροῦν,\FnParse{d}{present active participle nominative neuter singular} ὅτι τὸ πνεῦμά ἐστιν\FnParse{e}{present active indicative 3rd singular} ἡ ἀλήθεια. 
\MainTextVerseMark{5}{7}ὅτι τρεῖς εἰσιν\FnParse{a}{present active indicative 3rd plural} οἱ μαρτυροῦντες,\FnParse{b}{present active participle nominative masculine plural} 
\MainTextVerseMark{5}{8}τὸ πνεῦμα καὶ τὸ ὕδωρ καὶ τὸ αἷμα, καὶ οἱ τρεῖς εἰς τὸ ἕν εἰσιν.\FnParse{a}{present active indicative 3rd plural} 
\MainTextVerseMark{5}{9}εἰ τὴν μαρτυρίαν τῶν ἀνθρώπων λαμβάνομεν,\FnParse{a}{present active indicative 1st plural} ἡ μαρτυρία τοῦ θεοῦ μείζων ἐστίν,\FnParse{b}{present active indicative 3rd singular} ὅτι αὕτη ἐστὶν\FnParse{c}{present active indicative 3rd singular} ἡ μαρτυρία τοῦ θεοῦ ⸀ὅτι μεμαρτύρηκεν\FnParse{d}{perfect active indicative 3rd singular} περὶ τοῦ υἱοῦ αὐτοῦ. 
\MainTextVerseMark{5}{10}ὁ πιστεύων\FnParse{a}{present active participle nominative masculine singular} εἰς τὸν υἱὸν τοῦ θεοῦ ἔχει\FnParse{b}{present active indicative 3rd singular} τὴν μαρτυρίαν ἐν ⸀αὑτῷ· ὁ μὴ πιστεύων\FnParse{c}{present active participle nominative masculine singular} τῷ θεῷ ψεύστην\FnParseFormGloss{d}{accusative masculine singular}{ψεύστης}{liar} πεποίηκεν\FnParse{e}{perfect active indicative 3rd singular} αὐτόν, ὅτι οὐ πεπίστευκεν\FnParse{f}{perfect active indicative 3rd singular} εἰς τὴν μαρτυρίαν ἣν μεμαρτύρηκεν\FnParse{g}{perfect active indicative 3rd singular} ὁ θεὸς περὶ τοῦ υἱοῦ αὐτοῦ. 
\MainTextVerseMark{5}{11}καὶ αὕτη ἐστὶν\FnParse{a}{present active indicative 3rd singular} ἡ μαρτυρία, ὅτι ζωὴν αἰώνιον ἔδωκεν\FnParse{b}{aorist active indicative 3rd singular} ⸂ὁ θεὸς ἡμῖν⸃, καὶ αὕτη ἡ ζωὴ ἐν τῷ υἱῷ αὐτοῦ ἐστιν.\FnParse{c}{present active indicative 3rd singular} 
\MainTextVerseMark{5}{12}ὁ ἔχων\FnParse{a}{present active participle nominative masculine singular} τὸν υἱὸν ἔχει\FnParse{b}{present active indicative 3rd singular} τὴν ζωήν· ὁ μὴ ἔχων\FnParse{c}{present active participle nominative masculine singular} τὸν υἱὸν τοῦ θεοῦ τὴν ζωὴν οὐκ ἔχει.\FnParse{d}{present active indicative 3rd singular} 
\MainTextVerseMark{5}{13}Ταῦτα ἔγραψα\FnParse{a}{aorist active indicative 1st singular} ⸀ὑμῖν ἵνα εἰδῆτε\FnParse{b}{perfect active subjunctive 2nd plural} ὅτι ζωὴν ⸂ἔχετε\FnParse{c}{present active indicative 2nd plural} αἰώνιον⸃, ⸂τοῖς πιστεύουσιν⸃\FnParse{d}{present active participle dative masculine plural} εἰς τὸ ὄνομα τοῦ υἱοῦ τοῦ θεοῦ. 
\MainTextVerseMark{5}{14}καὶ αὕτη ἐστὶν\FnParse{a}{present active indicative 3rd singular} ἡ παρρησία ἣν ἔχομεν\FnParse{b}{present active indicative 1st plural} πρὸς αὐτόν, ὅτι ἐάν τι αἰτώμεθα\FnParse{c}{present middle subjunctive 1st plural} κατὰ τὸ θέλημα αὐτοῦ ἀκούει\FnParse{d}{present active indicative 3rd singular} ἡμῶν. 
\MainTextVerseMark{5}{15}καὶ ἐὰν οἴδαμεν\FnParse{a}{perfect active indicative 1st plural} ὅτι ἀκούει\FnParse{b}{present active indicative 3rd singular} ἡμῶν ὃ ⸀ἐὰν αἰτώμεθα,\FnParse{c}{present middle subjunctive 1st plural} οἴδαμεν\FnParse{d}{perfect active indicative 1st plural} ὅτι ἔχομεν\FnParse{e}{present active indicative 1st plural} τὰ αἰτήματα\FnParseFormGloss{f}{accusative neuter plural}{αἴτημα}{request} ἃ ᾐτήκαμεν\FnParse{g}{perfect active indicative 1st plural} ⸀ἀπ’ αὐτοῦ. 
\MainTextVerseMark{5}{16}ἐάν τις ἴδῃ\FnParse{a}{aorist active subjunctive 3rd singular} τὸν ἀδελφὸν αὐτοῦ ἁμαρτάνοντα\FnParse{b}{present active participle accusative masculine singular} ἁμαρτίαν μὴ πρὸς θάνατον, αἰτήσει,\FnParse{c}{future active indicative 3rd singular} καὶ δώσει\FnParse{d}{future active indicative 3rd singular} αὐτῷ ζωήν, τοῖς ἁμαρτάνουσιν\FnParse{e}{present active participle dative masculine plural} μὴ πρὸς θάνατον. ἔστιν\FnParse{f}{present active indicative 3rd singular} ἁμαρτία πρὸς θάνατον· οὐ περὶ ἐκείνης λέγω\FnParse{g}{present active indicative 1st singular} ἵνα ἐρωτήσῃ.\FnParse{h}{aorist active subjunctive 3rd singular} 
\MainTextVerseMark{5}{17}πᾶσα ἀδικία\FnParseFormGloss{a}{nominative feminine singular}{ἀδικία}{wickedness} ἁμαρτία ἐστίν,\FnParse{b}{present active indicative 3rd singular} καὶ ἔστιν\FnParse{c}{present active indicative 3rd singular} ἁμαρτία οὐ πρὸς θάνατον. 
\MainTextVerseMark{5}{18}Οἴδαμεν\FnParse{a}{perfect active indicative 1st plural} ὅτι πᾶς ὁ γεγεννημένος\FnParse{b}{perfect passive participle nominative masculine singular} ἐκ τοῦ θεοῦ οὐχ ἁμαρτάνει,\FnParse{c}{present active indicative 3rd singular} ἀλλ’ ὁ γεννηθεὶς\FnParse{d}{aorist passive participle nominative masculine singular} ἐκ τοῦ θεοῦ τηρεῖ\FnParse{e}{present active indicative 3rd singular} ⸀αὐτόν, καὶ ὁ πονηρὸς οὐχ ἅπτεται\FnParse{f}{present middle indicative 3rd singular} αὐτοῦ. 
\MainTextVerseMark{5}{19}οἴδαμεν\FnParse{a}{perfect active indicative 1st plural} ὅτι ἐκ τοῦ θεοῦ ἐσμεν,\FnParse{b}{present active indicative 1st plural} καὶ ὁ κόσμος ὅλος ἐν τῷ πονηρῷ κεῖται.\FnParseFormGloss{c}{present middle indicative 3rd singular}{κεῖμαι}{to lay} 
\MainTextVerseMark{5}{20}οἴδαμεν\FnParse{a}{perfect active indicative 1st plural} δὲ ὅτι ὁ υἱὸς τοῦ θεοῦ ἥκει,\FnParseFormGloss{b}{present active indicative 3rd singular}{ἥκω}{to come} καὶ δέδωκεν\FnParse{c}{perfect active indicative 3rd singular} ἡμῖν διάνοιαν\FnParseFormGloss{d}{accusative feminine singular}{διάνοια}{mind} ἵνα ⸀γινώσκωμεν\FnParse{e}{present active subjunctive 1st plural} τὸν ἀληθινόν·\FnParseFormGloss{f}{accusative masculine singular}{ἀληθινός}{true} καὶ ἐσμὲν\FnParse{g}{present active indicative 1st plural} ἐν τῷ ἀληθινῷ,\FnParseFormGloss{h}{dative masculine singular}{ἀληθινός}{true} ἐν τῷ υἱῷ αὐτοῦ Ἰησοῦ Χριστῷ. οὗτός ἐστιν\FnParse{i}{present active indicative 3rd singular} ὁ ἀληθινὸς\FnParseFormGloss{j}{nominative masculine singular}{ἀληθινός}{true} θεὸς καὶ ζωὴ αἰώνιος. 
\MainTextVerseMark{5}{21}Τεκνία,\FnParseFormGloss{a}{neuter plural}{τεκνίον}{dear children} φυλάξατε\FnParse{b}{aorist active imperative 2nd plural} ἑαυτὰ ἀπὸ τῶν ⸀εἰδώλων.\FnParseFormGloss{c}{genitive neuter plural}{εἴδωλον}{idol} 
\end{document}